\section{Technical Overview}
\label{sec:overview}

This section is a self-contained tour of the construction: a reader who
reads only the introduction and this overview will understand how the
argument works end to end, while the sections that follow supply the
definitions, lemmas, and proofs the tour forward-references. The spine has
four moves and one unifying property. The arithmetization makes every
constraint linear except a handful of column products
(\Cref{sec:overview-rep}); a binary-field polynomial IOP discharges the
linear constraints (\Cref{sec:overview-linear}); a single change of reading
lets the products ride that \emph{same} IOP (\Cref{sec:overview-crux}); and
a bit-sliced commitment settles everything with one opening
(\Cref{sec:overview-open}). What ties the four together is that every map
used to collapse a cell to the field is $\Ftwo$-linear
(\Cref{sec:overview-thread}).

\subsection{The representation: every primitive but $\AND$ is linear}
\label{sec:overview-rep}

Represent a $\wbit$-bit word as a bit-polynomial $\bp{u}=\sum_b u_b X^b$ in
the cell ring $\cellring=\Ftwo[X]/(X^{\wbit})$, the coefficient of $X^b$
being bit $b$ (\Cref{sec:arith}). In this encoding the bitwise primitives
of SHA-256 are ring operations rather than emulations:
\begin{itemize}[leftmargin=2em]
  \item \textbf{$\XOR$ is addition} ($1+1=0$). A column that is an
    $\Ftwo$-linear combination of others is never committed; the verifier
    forms the combination of the committed columns' projections when needed.
    The dozens of $\XOR$s in SHA-256 cost nothing.
  \item \textbf{Shifts and rotations are monomial multiplication} by $X^c$
    (in the shift quotient $\Fshift$ or rotation quotient $\Frot$), hence
    $\Ftwo$-linear; a shifted column is \emph{virtual} --- reconstructed
    from its source, not committed.
  \item \textbf{Modular addition is the Binius carry identity}, taken in the
    shift quotient $\Fshift=\Ftwo[X]/(X^{\wbit})$, where multiplication by $X$
    is a left shift and the relation $X^{\wbit}=0$ discards the top overflow ---
    exactly the $\bmod\,2^{\wbit}$ wrap. For a binary step
    $\bp{s}\equiv\bp{a}+\bp{b}\pmod{2^{\wbit}}$ the carry word $\bp{c}$ is the
    \emph{down-shift} of the free combination $\bp{s}+\bp{a}+\bp{b}$, namely
    $\bp{c}[i]=(\bp{s}+\bp{a}+\bp{b})[i+1]$. Its lower $\wbit-1$ bits are thus
    the virtual column $\SHR^{1}(\bp{s}+\bp{a}+\bp{b})$, reconstructed for free,
    while its top bit $\bp{c}[\wbit-1]=\beta$ --- the carry-out of the
    most-significant position, lost to the wrap and \emph{not} an
    $\Ftwo$-linear function of $\bp{s},\bp{a},\bp{b}$ --- is the single
    committed bit. With $\bp{c}$ so defined the step is one linear identity ---
    membership of $\bp{s}+\bp{a}+\bp{b}$ in $\ideal{X}$, whose only content is
    the bit-$0$ equation $(\bp{s}+\bp{a}+\bp{b})[0]=0$ --- together with one
    full-adder Hadamard, both relations living in $\Fshift$:
    \[
      \bp{s}=\bp{a}+\bp{b}+X\bp{c},
      \qquad
      (\bp{a}+X\bp{c})\had(\bp{b}+X\bp{c}) = \bp{c}+X\bp{c}.
    \]
    Multiplying the carry by $X$ shifts it back up: $X\bp{c}$ restores bits
    $1,\dots,\wbit-1$ of the sum and sends $\beta$ into the discarded
    $X^{\wbit}$ term, so the lone committed bit re-enters exactly as the
    overflow it represents (\Cref{lem:binius}).
\end{itemize}

Committing a \emph{single} bit per addition is a contribution we highlight,
and a direct dividend of the cell-ring representation: the carry word is, but
for its overflow bit, a shift of committed data, and a shift is a free virtual
column here (first item above), so the prover commits only the one carry-out
bit $\beta$ that the shift cannot supply --- never a carry \emph{word}. The
contrast with \binius{} is the same shift mechanism seen from the witness side.
\binius{} uses the identical carry identity, but over packed words whose shifts
are \emph{not} free read-offs (they are discharged by a dedicated Shift
reduction, \Cref{sec:shifts}); its adder gadget therefore materialises a full
$\wbit$-bit carry word $\mathsf{cout}$ as a committed wire, constrained by its
own full-adder identity
\[
  \bigl(\bp{a}\oplus(\mathsf{cout}\ll 1)\bigr)\wedge
  \bigl(\bp{b}\oplus(\mathsf{cout}\ll 1)\bigr)
  \;=\; \mathsf{cout}\oplus(\mathsf{cout}\ll 1),
\]
the shift $\mathsf{cout}\ll 1$ playing the role of our $X\bp{c}$. Where
\binius{} commits $\wbit$ carry bits per addition, the all-$\Ftwo[X]$
arithmetization commits one.
The lone nonlinear primitive is bitwise $\AND$, the coefficient-wise
\emph{Hadamard} product $\bp{u}\had\bp{v}$ (\Cref{lem:and}); $\Ch$ and
$\Maj$ each rewrite to a single such product. The arithmetization therefore
produces an AIR in which \emph{every constraint is a linear
ideal-membership relation except a small fixed number of column Hadamard
products} --- in the deployed layout $K_H=14$ of them ($2$ for the $\AND$s,
$12$ for the addition carries). The rest of the proof has exactly two jobs:
the linear constraints, and the Hadamard products.

\subsection{The linear part: ideal check, the projection $\psia$, sumcheck}
\label{sec:overview-linear}

The prover commits the primary witness columns natively in $\cellring$ with
a hash-based commitment whose $\Ftwo$-linear code \emph{does not widen} the
cell ring (\Cref{sec:commitment}). Because the trace already lives in
$\Ftwo[X]\subset\K[X]$, the prime-field projection that opens the general
$\zinc$ compiler is vacuous --- the first concrete saving of working in
characteristic two.

The linear constraints are ideal-membership assertions
$Q_\ell(T)\in\ideal{g_\ell}$ with generators
$g_\ell\in\{0,\,X\}$ (equality/boundary and the addition LSB); rotations
and shifts never appear as constraints --- they are virtual columns --- so
the ideal check stays entirely within $\cellring=\Ftwo[X]/(X^{\wbit})$.
The \emph{ideal check} (\Cref{sec:ideal-check}) batches all of
them --- across families and across all rows --- into combined polynomials
$E_\ell\in\K[X]$ carrying the obligation $E_\ell\in\ideal{g_\ell}$. The
\emph{evaluation projection} $\psia$ then substitutes the cell variable
$X=\alpha$ for a transcript-fresh $\alpha\in\K$. Linear constraints
\emph{survive} this projection: since $\psia$ is $\Ftwo$-linear and
evaluation is multiplicative, \albert{THe right arrow below is not great/is confusing. We need the two directions, but reverse implication only holds with overwhelming probability due to Zinc+ work. We should be more clear about it}
\[
  \psia\bigl(Q(T)\bigr)=Q\bigl(\psia(T)\bigr),
  \qquad
  Q(T)\in\ideal{g}\ \Longrightarrow\ \psia\bigl(Q(T)\bigr)=g(\alpha)\cdot u
\]
(\Cref{lem:linear-survives}), so a linear identity among cells becomes a
linear identity among the projected columns, and membership becomes a plain
$\K$-identity once $\alpha$ is fixed. A standard degree-$2$ sumcheck
(\Cref{sec:sumcheck}) reduces the resulting zerocheck over the $2^{\nu}$
rows to per-column evaluation claims $\mle{\psia(w_g)}(r^*)$ at a single
point $r^*$.

\subsection{The crux: $\AND$ rides the same pipeline through a second reading $\psiz$}
\label{sec:overview-crux}

The Hadamard products cannot ride $\psia$, because evaluation does not
respect the coefficient-wise product:
\[
  \psia(\bp{u}\had\bp{v}) = \sum_b u_b v_b\,\alpha^{b}
  \;\neq\;
  \psia(\bp{u})\,\psia(\bp{v}) = \sum_{b,b'} u_b v_{b'}\,\alpha^{b+b'},
\]
the cross-terms $b\neq b'$ scrambling the bit-products
(\Cref{rem:and-not-survive}). The resolution is to read the \emph{same}
committed bits under a \emph{different} weighting. Identify the $\wbit$ bit
positions with an $\Ftwo$-affine subspace $H\subset\K$ and read a cell's
bits as the \emph{Lagrange} coefficients $L(\bp{u})=\sum_b u_b L_b$ of a
univariate over $H$, rather than as monomial coefficients. Then
\[
  \bp{u}\had\bp{v}=\bp{w}
  \quad\Longleftrightarrow\quad
  L(\bp{u})\,L(\bp{v})-L(\bp{w}) \in \ideal{Z_H},
  \qquad Z_H=\textstyle\prod_{h\in H}(X-h)
\]
(\Cref{eq:had-ideal}): the product is correct on every bit exactly when
this polynomial vanishes on $H$. That is an \emph{ideal-membership claim}
--- the very shape the linear part already discharges, now with the extra
generator $Z_H$. So the nonlinearity rides the same machinery: the same
ideal check (the prover ships only the $\wbit-1$ off-subspace values of the
combined product polynomial, which vanish on $H$ for honest data), the same
point $\alpha$ --- now reading the cell in the \emph{Lagrange} basis,
\[
  \psiz(\bp{w}) = L(\bp{w})(\alpha) = \sum_b L_b(\alpha)\,w_b,
\]
rather than the monomial one --- and one extra degree-$3$ group in the
\emph{same} sumcheck, reduced to the \emph{same} point $r^*$
(\Cref{sec:hadamard}). The Lagrange reading keeps the bit-products $u_bv_b$
separated by position, so the verifier's closing check is a genuine
per-position product rather than the scrambled cross-terms
$\psia(\bp{u}\had\bp{v})$ would give. $\psia$ and $\psiz$ are not two
projections: they are the one substitution $X=\alpha$ with the bits read in
the monomial resp.\ Lagrange basis (the weights $\alpha^b=X^b|_\alpha$ and
$L_b(\alpha)=L_b(X)|_\alpha$ are the same point in two bases). The
$\mu=\log_2\wbit$ bit positions are folded once, by this Lagrange reading,
and never become sumcheck variables.

\subsection{One opening answers both readings}
\label{sec:overview-open}

After the sumcheck and a multipoint-evaluation step (\Cref{sec:multipoint}),
both halves have reduced to multilinear-evaluation claims at one common
point $r_0$: the linear claims are $\psia$-images of the committed columns,
the discharge claims their $\psiz$-images. The commitment settles both from
a single object. Its opening binds, per column, the \emph{bit-slice fold}
\[
  a'_g=\sum_{b}a'_{g,b}\,X^b\in\K[X]_{<\wbit},
  \qquad
  a'_{g,b}=\sum_{x\in\{0,1\}^{\nu}}\eq_x(r_0)\,(w_g)_b(x),
\]
a degree-$<\wbit$ polynomial over $\K$ whose $b$-th coefficient is the
$b$-th bit slice of the column, folded along $\eq(\cdot;r_0)$. Because every
cell projection is the inner product of the coefficient vector with its
weights, \emph{any} projected evaluation at $r_0$ is a single
length-$\wbit$ read-off of $a'_g$:
\[
  \psia(a'_g)=\mle{\psia(w_g)}(r_0)\ \ (\text{linear claim}),
  \qquad
  \psiz(a'_g)=\mle{\psiz(w_g)}(r_0)\ \ (\text{discharge claim}).
\]
One bound object answers both readings, so the nonlinear constraints cost
\emph{no second opening and no second point} (\Cref{sec:commitment}). The
opening itself is a flat Brakedown/Ligero-style proximity test over the
$\Ftwo$-RAA code, with knowledge-soundness error
$\epsilon_{\mathrm{prox}}(t,\delta)+O(1/|\K|)$ dominated by its $t$ column
queries.

\subsection{Provenance and the proof-size lever: Blaze and code-switching}
\label{sec:overview-blaze}

The $\Ftwo$-RAA commitment used above
(\Cref{sec:overview-open,sec:commitment}) sits in the family introduced by
\emph{Blaze} \cite{blaze}, a hash-based multilinear polynomial commitment
over binary fields. Blaze's commit is dominated by the XOR-only encoding of
an interleaved Repeat--Accumulate--Accumulate code and a single Merkle tree
--- the same linear-time, non-widening code our opening rests on --- which
it then opens not by a flat proximity test but by \emph{code-switching}: the
evaluation/proximity claim on the RAA codeword is reduced to a claim on a
Reed--Solomon codeword and discharged by an inner foldable IOP of proximity
(BaseFold/WHIR) \cite{basefold,whir}, yielding a polylogarithmic verifier.
Our construction shares Blaze's code and its linear-time prover; what it adds
is the \emph{bit-sliced} opening (\Cref{sec:overview-open}): Blaze, like the
IOPPs it composes, settles a single evaluation, whereas binding the
width-$\wbit$ fold $a'$ lets the two $\Ftwo$-linear readings $\psia$ and
$\psiz$ share \emph{one} opening, so the Hadamard constraints ride the linear
part at no extra commitment cost. That dual read-off is the feature this
paper contributes on top of the Blaze code; it has no analogue in a
single-evaluation MLPCS.

The two designs differ only in how the opening is \emph{settled}. We deploy
the \emph{flat} Brakedown/Ligero proximity test \cite{brakedown}, whose $t$
column queries dominate proof size (\Cref{sec:overview-open}) and account for
the bulk of the $4$--$7\times$ proof-size gap to Binius64's recursive scheme
(\Cref{sec:cmp}); Blaze's code-switching is exactly the route from that flat
test to a recursive one over the \emph{same} fast code --- the proof-size
lever already sketched in \Cref{sec:cmp-outlook}. Crucially, this lever is
compatible with the bit-sliced reading. The fold axis (the $\nu$ row
variables, weighted by $\eq(\cdot;r_0)$) is orthogonal to the cell axis (the
$\wbit$ bit positions of $X$), and the $\eq$-fold, the code, and every
BaseFold/WHIR round operation are $\Ftwo$- (resp.\ $\K$-) linear
\emph{coordinate-wise on the cell axis}. A recursive opening therefore still
certifies the full width-$\wbit$ object $a'\in\Kpoly$ rather than a single
collapsed value, provided the cell axis is carried as passive shared width
--- never folded, batched, or projected inside the recursion --- so that
$\psia(a')$ and $\psiz(a')$ are read off afterward exactly as in
\Cref{sec:overview-open}. The one-opening-answers-both property is thus
independent of whether the proximity test is flat or recursive: a
code-switched opening would shrink the proof toward the recursive class and
give an $O(\polylog)$ verifier while preserving the $\psia/\psiz$ sharing.
The deployed system uses the flat test; this is recorded as the concrete
direction for closing the proof-size gap.

\subsection{The thread, the guarantee, and the scope}
\label{sec:overview-thread}

One property runs through all four moves: every projection $\psi_\xi$ is an
$\Ftwo$-linear map $\cellring\to\K$, $\sum_b w_b X^b\mapsto\sum_b
w_b\,\xi_b$, and the commitment's code is $\Ftwo$-linear and does not widen
$\cellring$. Hence projections commute through encoding, $\XOR$ of columns
commutes with projection, shift and rotation virtuals commute with both,
the random projection may be applied \emph{after} committing, and one
committed object answers many functionals (\Cref{rem:linearity-thread}).
That single fact is what makes $\XOR$ free, shifts virtual, and the
$\psia$/$\psiz$ pair share an opening.

\begin{remark}[Why the entry-wise shift is a free read-off, and why a packed representation is not]
\label{rem:shift-readoff}
It is worth isolating \emph{why} a shift costs nothing here, because it is the
property a competing representation most easily loses. A shift is an
$\Ftwo$-linear map on a cell's $\wbit$ bits, and the opening binds the
\emph{entire} width-$\wbit$ fold $a'_g$ (\Cref{sec:overview-open}), not a single
collapsed value. The projection of a shifted column is therefore one more
length-$\wbit$ read-off of that \emph{same} object: in the truncating
convention of \Cref{lem:shr},
\[
  \psia(\SHR^{r}\bp{u})
  \;=\; \sum_{b=0}^{\wbit-1-r} u_{b+r}\,\alpha^{b}
  \;=\; \alpha^{-r}\,\psia(\bp{u}) \;+\; \sum_{b<r} u_b\,\alpha^{\,b-r}
  \qquad(\text{over }\Ftwo),
\]
manifestly a length-$\wbit$ functional of the cell's bits (the first form), and
equivalently a fixed $\alpha^{-r}$ twist of the column's own projection plus a
fixed $\Ftwo$-linear functional of the few bits that fall off the boundary (the
second); both are read off $a'_g$ and absorbed into the $O(\log N)$ multipoint
step (\Cref{sec:multipoint}). A rotation is the same with the wrapped bits
re-entering rather than vanishing (\Cref{lem:rot}). No shifted column is
committed and no dedicated argument runs. The boundary term is precisely the
obstruction to reading the shift off as the bare scalar $\alpha^{-r}\psia(\bp{u})$:
$\psia$ is an algebra homomorphism on $\Ftwo[X]$, but the cell ring's relation
$X^{\wbit}=0$ is not respected by evaluation at $\alpha$, so the wrapped bits
survive as that correction. The enabling fact is structural and is exactly the
thread above --- the cell axis (the $\wbit$ positions of $X$) is passive shared
width \emph{orthogonal to the code}, which does not widen $\cellring$, so any
$\Ftwo$-linear map on that axis is answered by reading one more functional off
the bound object.

A representation that instead places the bit positions on the commitment's
\emph{encoded} axis cannot do this, and this is \binius{}'s position
(\Cref{sec:shifts}): a word's bits are hypercube coordinates of the committed
multilinear, encoded and folded by the Reed--Solomon/FRI commitment, so a shift
permutes the very coordinates the code acts on. The shifted operand's evaluation
is then a contraction $\sum_{i,j} g(i,j)\,h(i,j)$ pairing the \emph{secret}
witness $g$ with the public shift indicator $h$, which the verifier cannot form
unaided and which a dedicated \emph{Shift reduction} sumcheck discharges ---
\binius{}'s single largest prover phase. The distinction is sharp but narrow,
and two qualifications keep it honest. \binius{}'s shifted operands are
themselves \emph{virtual}: no shifted copy is committed, and the entry-wise
shift is succinct in $N$ in both systems. What differs is only the
\emph{prover-side cost} of tying the shift to the committed word --- a free
read-off here, a sumcheck there --- and that is a property of the chosen
\emph{representation}, not a barrier intrinsic to binary fields: a system that
carried the bit axis as an internal algebra dimension would recover the free
shift, converging toward the cell-ring AIR. It is moreover the cross-word
\emph{wiring}, not the entry-wise shift, that makes \binius{}'s \emph{verifier}
linear (\Cref{obs:wiring-not-shifts}); the shift gap is purely a prover-side one.
\end{remark}

Composed under Fiat--Shamir, the argument is a knowledge-sound SNARK for the
SHA-256 chained-compression relation, with error
\[
  \epsilon \;\le\; \epsilon_{\mathrm{prox}}(t,\delta)
  + O\!\Bigl(\tfrac{N_{\mathrm{con}}\,\wbit + K_H + \nu}{|\K|}\Bigr)
\]
(\Cref{thm:e2e}), the proximity term dominating. Two honest caveats remain
(\Cref{sec:caveats}): the construction is not zero-knowledge, and the
modular-addition carries are presently bound by a trusted reading rather
than to the commitment --- the $\AND$ relations and the entire linear part
are unconditional, and closing the adder gap is a localized, recorded
follow-up. The final part of the paper (\Cref{sec:cmp}) places the
construction against Binius64, the closest system: on one machine at a
matched commitment rate the $\Ftwo[X]$ prover leads and the lead widens with
input size, its verifier runs from near-parity to an order of magnitude
faster ($\sqrt N$ versus $O(N)$), while Binius64's recursive proof is
several times smaller.
